% Part: many-valued-logic
% Chapter: three-valued-logics
% Section: multiple-designation

\documentclass[../../../include/open-logic-section]{subfiles}

\begin{document}

\olfileid{mvl}{thr}{mul}

\olsection{শুধু~$\True$ নয়, অন্য মানও মনোনীত করা}

এ পর্যন্ত দেখা সব যুক্তিবিদ্যায় মনোনীত সত্যমানের সেট ছিল
$V^+ = \{\True\}$; অর্থাৎ কোনো কিছুর সত্যমান~$\True$ হলেই এবং কেবল
তখনই সেটি সত্য বলে গণ্য হতো। কিন্তু কোনো কিছুর সত্যমান শুধু~$\False$
না হলেই তাকেও সত্য গণ্য করা যেতে পারে। তখন ম্যাট্রিক্সে, যেমন,
$V^+ = \{\True, \Undef\}$ স্থির করে একটি যুক্তিবিদ্যা পাওয়া যাবে।

\begin{defn}
\emph{কূটাভাসের যুক্তিবিদ্যা}~$\LogLP$ নিচের ম্যাট্রিক্স দিয়ে সংজ্ঞায়িত:
\begin{enumerate}
  \item $\lnot$, $\land$, $\lor$, $\lif$-সহ প্রমিত বচনগত ভাষা
  $\Lang L_0$।
  \item সত্যমানের সেট $V = \{\True, \Undef, \False\}$।
  \item $\True$ ও~$\Undef$ মনোনীত; অর্থাৎ
  $V^+ = \{\True, \Undef\}$।
  \item সত্যমান-অপেক্ষকগুলি শক্তিশালী ক্লিনি যুক্তিবিদ্যার মতোই।
\end{enumerate}
\end{defn}

\begin{defn}
Halld\'en-এর \emph{অর্থহীনতার যুক্তিবিদ্যা}~$\LogHal$ নিচের ম্যাট্রিক্স
দিয়ে সংজ্ঞায়িত:
\begin{enumerate}
  \item $\lnot$, $\land$, $\lor$, $\lif$ এবং $1$-স্থানীয় সংযোজক~$+$-সহ
  প্রমিত বচনগত ভাষা $\Lang L_0$।
  \item সত্যমানের সেট $V = \{\True, \Undef, \False\}$।
  \item $\True$ ও~$\Undef$ মনোনীত; অর্থাৎ
  $V^+ = \{\True, \Undef\}$।
  \item সত্যমান-অপেক্ষকগুলি দুর্বল ক্লিনি যুক্তিবিদ্যার মতো; সঙ্গে আছে
  “অর্থহীন” অপারেটর:
  \begin{center}
    \begin{tabular}{c|c}
    $\tf{+}$ & \\
    \hline
    $\True$ & $\False$ \\
    $\Undef$ & $\True$ \\
    $\False$ & $\False$
    \end{tabular}
  \end{center}
\end{enumerate}
\end{defn}

যে ক্লিনি যুক্তিবিদ্যাগুলির সঙ্গে এদের সত্যসারণি অভিন্ন, তাদের বিপরীতে
এই যুক্তিবিদ্যাগুলিতে সর্বতঃসত্য \emph{আছে}।

\begin{prop}\ollabel{prop:LP-taut-CL} $\LogLP$-এর সর্বতঃসত্যগুলি ধ্রুপদি
  বচনগত যুক্তিবিদ্যার সর্বতঃসত্যগুলির সঙ্গে অভিন্ন।
\end{prop}

\begin{proof}
  \olref[syn][sub]{prop:mvl-cl} অনুযায়ী,
  $\Entails[\LogLP] !A$ হলে $\Entails[\LogCL] !A$। বিপরীত দিকটি দেখাতে
  আমরা প্রমাণ করব: যদি এমন !!a{valuation}
  $\pAssign v\colon \PVar \to \{\False, \True, \Undef\}$ থাকে যাতে
  $\pValue v(!A)[\LogKs] = \False$, তবে এমন !!a{valuation}
  $\pAssign {v'}\colon \PVar \to \{\False, \True\}$ আছে যাতে
  $\pValue {v'}(!A)[\LogCL] = \False$। এটি $\LogLP$-এর জন্য ফলটি
  প্রতিষ্ঠা করে; কারণ $\LogKs$ ও~$\LogLP$-এর বৈশিষ্ট্যসূচক
  সত্যমান-অপেক্ষক অভিন্ন, আর $\LogLP$-এ $\False$-ই একমাত্র অমনোনীত
  সত্যমান ($\LogLP$ ও~$\LogKs$-এর মধ্যে এটিই একমাত্র পার্থক্য)। অতএব
  $\Entails/[\LogLP] !A$ হলে কোনো !!{valuation}~$\pAssign v$-এর জন্য
  $\pValue v(!A)[\LogLP] = \pValue v[\LogKs](!A) = \False$। আমরা যে
  দাবিটি প্রমাণ করছি, তা থেকে
  $\pValue{v'}[\LogCL](!A) = \False$, অর্থাৎ
  $\Entails/[\LogCL] !A$।

  দাবিটি প্রতিষ্ঠা করতে প্রথমে $\pAssign {v'}$-কে এভাবে সংজ্ঞায়িত করি:
  \[
    \pAssign {v'}(p) =
    \begin{cases}
    \True & \text{যদি } \pAssign {v}(p) \in \{\True, \Undef\}\\
    \False & \text{অন্যথায়।}
  \end{cases}
  \]
  এখন $!A$-এর উপর আরোহের সাহায্যে দেখাই যে (ক)~যদি
  $\pValue v(!A)[\LogKs] = \False$ হয়, তবে
  $\pValue {v'}(!A)[\LogCL] = \False$; এবং (খ)~যদি
  $\pValue v(!A)[\LogKs] = \True$ হয়, তবে
  $\pValue {v'}(!A)[\LogCL] = \True$।
  \begin{enumerate}
    \item আরোহের ভিত্তি: $!A \ident p$।
    \olref[syn][val]{defn:pValue} ও~$v'$-এর সংজ্ঞা অনুযায়ী, যদি
    $\pValue v(!A)[\LogKs] = \False$ হয়, তবে
    $\pAssign v(p)=\False$; উপরের সংজ্ঞা থেকে
    $\pValue {v'}(!A)[\LogCL]=\False$। আবার যদি
    $\pValue v(!A)[\LogKs] = \True$ হয়, তবে
    $\pAssign v(p)=\True$; তাই
    $\pValue {v'}(!A)[\LogCL]=\True$। ফলে (ক) ও~(খ) উভয়ই সিদ্ধ।

    আরোহ-ধাপের জন্য ক্ষেত্রগুলি বিবেচনা করি:
    \item $!A \ident \lnot !B$।
    \begin{enumerate}
      \item ধরা যাক, $\pValue v(\lnot!B)[\LogKs] = \False$।
      $\tf{\lnot}[\LogKs]$-এর সংজ্ঞা অনুযায়ী
      $\pValue v(!B)[\LogKs] = \True$। আরোহানুমান, ক্ষেত্র~(খ), থেকে
      $\pValue {v'}(!B)[\LogCL] = \True$; তাই
      $\pValue {v'}(\lnot !B)[\LogCL] = \False$।
      \item ধরা যাক, $\pValue v(\lnot!B)[\LogKs] = \True$।
      $\tf{\lnot}[\LogKs]$-এর সংজ্ঞা অনুযায়ী
      $\pValue v(!B)[\LogKs] = \False$। আরোহানুমান, ক্ষেত্র~(ক), থেকে
      $\pValue {v'}(!B)[\LogCL] = \False$; তাই
      $\pValue {v'}(\lnot !B)[\LogCL] = \True$।
    \end{enumerate}

    \item $!A \ident (!B \land !C)$।
    \begin{enumerate}
      \item ধরা যাক, $\pValue v(!B \land !C)[\LogKs] = \False$।
      $\tf{\land}[\LogKs]$-এর সংজ্ঞা অনুযায়ী
      $\pValue v(!B)[\LogKs] = \False$ অথবা
      $\pValue v(!C)[\LogKs] = \False$। আরোহানুমান, ক্ষেত্র~(ক), থেকে
      $\pValue {v'}(!B)[\LogCL] = \False$ অথবা
      $\pValue {v'}(!C)[\LogCL] = \False$; তাই
      $\pValue {v'}(!B \land !C)[\LogCL] = \False$।
      \item ধরা যাক, $\pValue v(!B \land !C)[\LogKs] = \True$।
      $\tf{\land}[\LogKs]$-এর সংজ্ঞা অনুযায়ী
      $\pValue v(!B)[\LogKs] = \True$ এবং
      $\pValue v(!C)[\LogKs] = \True$। আরোহানুমান, ক্ষেত্র~(খ), থেকে
      $\pValue {v'}(!B)[\LogCL] = \True$ এবং
      $\pValue {v'}(!C)[\LogCL] = \True$; তাই
      $\pValue {v'}(!B \land !C)[\LogCL] = \True$।
    \end{enumerate}
  \end{enumerate}

  অন্য দুটি ক্ষেত্র একই রকম; সেগুলি অনুশীলন হিসেবে রাখা হলো। বিকল্পভাবে,
  উপরের প্রমাণটি শুধু $\lnot$ ও~$\land$ থাকা সব !!{formula}s-এর জন্য ফলটি
  প্রতিষ্ঠা করে। এখন ব্যবহার করা যায় যে $\LogKs$ ও~$\LogCL$ উভয়েই,
  যেকোনো $\pAssign v$-এর জন্য,
  $\pValue v(!B \lor !C) = \pValue v(\lnot(\lnot!B \land \lnot !C))$
  এবং
  $\pValue v(!B \lif !C) = \pValue v(\lnot(!B \land \lnot !C))$।
\end{proof}

\begin{prob}
  \olref[mvl][thr][mul]{prop:LP-taut-CL}-এর প্রমাণ সম্পূর্ণ করো; অর্থাৎ
  $!A \ident (!B \lor !C)$ ও~$!A \ident (!B \lif !C)$ ক্ষেত্রগুলির
  জন্য (ক) ও~(খ) প্রতিষ্ঠা করো।
\end{prob}

\begin{prob}
প্রমাণ করো যে প্রতিটি ধ্রুপদি সর্বতঃসত্য~$\LogHal$-এ সর্বতঃসত্য।
\end{prob}

ধ্রুপদি যুক্তিবিদ্যার সঙ্গে তাদের সর্বতঃসত্যগুলি অভিন্ন হলেও তাদের
অনুসিদ্ধান্ত সম্পর্ক আলাদা। যেমন, $\LogLP$ \emph{পরাসঙ্গত}: এতে
$\lnot p, p \Entails/ q$; তাই বিস্ফোরণ নীতি
$\lnot !A, !A \Entails !B$ সাধারণভাবে সিদ্ধ নয়। ($!A$ ও~$!B$-এর কোনো
কোনো ক্ষেত্রে এটি সিদ্ধ; যেমন, $!B$ সর্বতঃসত্য হলে।)

\begin{prob}
  নিচের কোন সম্পর্কগুলি (ক)~$\LogLP$ এবং (খ)~$\LogHal$-এ সিদ্ধ?
  প্রতিটির জন্য একটি সত্যসারণি দাও।
  \begin{enumerate}
    \item $p, p \lif q \Entails q$
    \item $\lnot q, p \lif q \Entails \lnot p$
    \item $p \lor q, \lnot p \Entails q$
    \item $\lnot p, p \Entails q$
    \item $p \Entails p \lor q$
    \item $p \lif q, q\lif r \Entails p \lif r$
  \end{enumerate}
\end{prob}

$\LogLuk[3]$-এ $\Undef$-কেও মনোনীত করলে কী হয়?

\begin{defn}
  \emph{3-মানী R-Mingle যুক্তিবিদ্যা}~$\LogRM[3]$ নিচের ম্যাট্রিক্স
  দিয়ে সংজ্ঞায়িত:
  \begin{enumerate}
    \item $\lfalse$, $\lnot$, $\land$, $\lor$, $\lif$-সহ প্রমিত বচনগত
    ভাষা $\Lang L_0$।
    \item সত্যমানের সেট $V = \{\True, \Undef, \False\}$।
    \item $\True$ ও~$\Undef$ মনোনীত; অর্থাৎ
    $V^+ = \{\True, \Undef\}$।
    \item $\tf{\lfalse}=\False$; অবশিষ্ট সত্যমান-অপেক্ষকগুলি
    \L ukasiewicz যুক্তিবিদ্যা~$\LogLuk[3]$-এর মতোই।
  \end{enumerate}
\end{defn}

\begin{prob}
  নিচের কোন সম্পর্কগুলি $\LogRM[3]$-এ সিদ্ধ?
  \begin{enumerate}
    \item $p, p \lif q \Entails q$
    \item $p \lor q, \lnot p \Entails q$
    \item $\lnot p, p \Entails q$
    \item $p \Entails p \lor q$
  \end{enumerate}
\end{prob}

মনোনীত মানগুলি বদলালেই কখনও কখনও আলাদা সত্যসারণি থেকে একই যুক্তিবিদ্যা
(অনুসিদ্ধান্ত সম্পর্ক) পাওয়া যায়। যেমন, G\"odel যুক্তিবিদ্যায়
$V^+ = \{\True, \Undef\}$-কে~$\{\True\}$-এর বদলে নিলে এমনটি ঘটে।

\begin{prop}\ollabel{prop:gl-udes}
  $V = \{\False, \Undef,\True\}$, $V^+=\{\True, \Undef\}$ এবং
  $3$-মানী G\"odel যুক্তিবিদ্যার সত্যমান-অপেক্ষকবিশিষ্ট ম্যাট্রিক্সটি
  ধ্রুপদি যুক্তিবিদ্যা সংজ্ঞায়িত করে।
\end{prop}

\begin{proof}
  অনুশীলন।
\end{proof}

\begin{prob}
  \olref[mvl][thr][mul]{prop:gl-udes} প্রমাণ করো। এজন্য G\"odel
  যুক্তিবিদ্যার মতোই, কিন্তু $V^+=\{\True,\Undef\}$ নিয়ে সংজ্ঞায়িত
  যুক্তিবিদ্যা~$\Log L$-এর জন্য দেখাও যে
  $\Gamma \Entails/[\Log L] !B$ হলে
  $\Gamma \Entails/[\LogCL] !B$।
  \olref[mvl][thr][mul]{prop:LP-taut-CL}-এর ধারণা ব্যবহার করো; তবে (ক) ও
  (খ) বৈশিষ্ট্য প্রমাণের বদলে দেখাও যে
  $\pValue v(!A)[\LogGod] = \False$ ঠিক তখনই, যখন
  $\pValue {v'}(!A)[\LogCL] = \False$ (এবং তাই
  $\pValue v(!A)[\LogGod] \in \{\True,\Undef\}$ ঠিক তখনই, যখন
  $\pValue {v'}(!A)[\LogCL] = \True$)। এটি কেন প্রস্তাবটি প্রতিষ্ঠা
  করে তা ব্যাখ্যা করো।
\end{prob}
% BN-SRC-320 TARGET-CORRECTION-BEGIN
\par\smallskip\noindent\textnormal{\small\textbf{উৎস-সংশোধন (BN-SRC-320):} হিমায়িত LP প্রমাণের আরোহ-ভিত্তিতে $\pValue v(p)[\LogKs]=\pAssign v(p)=\pValue{v'}(p)[\LogCL]$ নিঃশর্তভাবে দাবি করা হয়েছে; কিন্তু $\pAssign v(p)=\Undef$ হলে সংজ্ঞা অনুযায়ী $\pAssign{v'}(p)=\True$, তাই সমতাটি মিথ্যা। বাংলা প্রমাণে দাবির জন্য প্রয়োজনীয় $\False$ ও~$\True$ ক্ষেত্র দুটি আলাদাভাবে দেখানো হয়েছে।} % BN-SRC-320 TARGET-CORRECTION-END
% BN-SRC-321 TARGET-CORRECTION-BEGIN
\par\smallskip\noindent\textnormal{\small\textbf{উৎস-সংশোধন (BN-SRC-321):} হিমায়িত LP প্রমাণের সংযোজন-ক্ষেত্রে $\pValue v(!B)$-কে “অথবা” ও “এবং” উভয় শাখায় ভুল করে দুবার লেখা হয়েছে। $!B\land !C$-এর শক্তিশালী ক্লিনি সত্যসারণি ও পরের ধ্রুপদি সিদ্ধান্ত অনুযায়ী দ্বিতীয় ঘটনাগুলিতে $\pValue v(!C)$ পুনরুদ্ধার করা হয়েছে।} % BN-SRC-321 TARGET-CORRECTION-END
% BN-SRC-322 TARGET-CORRECTION-BEGIN
\par\smallskip\noindent\textnormal{\small\textbf{উৎস-সংশোধন (BN-SRC-322):} হিমায়িত R-Mingle সংজ্ঞা ভাষায় শূন্য-স্থানীয় সংযোজক~$\lfalse$ রাখলেও তার সত্যমান-অপেক্ষক নির্ধারণ করেনি; “$\LogLuk[3]$-এর মতো” কথাটি এটি পূরণ করে না, কারণ পূর্ববর্তী $\LogLuk[3]$ ভাষায়~$\lfalse$ নেই। ঘোষিত ম্যাট্রিক্স সম্পূর্ণ করতে বাংলা লক্ষ্যে $\tf{\lfalse}=\False$ স্পষ্ট করা হয়েছে।} % BN-SRC-322 TARGET-CORRECTION-END

\end{document}
